%=============================================================================
% FlowClass.jl — Application Note for Bioinformatics (Oxford University Press)
%
% This file targets the OUP authoring class "oup-authoring-template.cls" as
% shipped with TeX Live (Document Class 2022/03/18 v1.1). It also compiles
% unchanged on Overleaf ("Oxford University Press - Contemporary").
%
% Note: the contemporary/modern header normally includes a society logo file
% (PNAS-Nexus-Logo-CMYK) that ships only with the Overleaf project, not the
% TeX Live package, so we blank \societylogo below to compile locally. Remove
% that line if you compile on Overleaf and want the journal furniture.
%
% Article type: Application Note (max ~2,600 words, or ~2,000 words + 1 figure).
% Bibliography uses paper.bib; style oup-abbrvnat.bst ships with the template.
%=============================================================================

\documentclass[unnumsec,webpdf,contemporary,large,namedate]{oup-authoring-template}

\usepackage{amsmath}
\usepackage{amssymb}
\usepackage{bm}
\usepackage{graphicx}

% Blank the society logo so the document compiles without the OUP logo file.
\makeatletter
\def\societylogo{}
\makeatother

% ---- Journal metadata -------------------------------------------------------
\journaltitle{Bioinformatics}
\DOI{DOI added during production}
\copyrightyear{2026}
\pubyear{2026}
\vol{00}
\issue{0}
\access{Advance Access Publication Date: }
\appnotes{Application Note}

%=============================================================================
\begin{document}

\title[FlowClass.jl]{FlowClass.jl: structural classification of dynamical
systems for landscape models of cell fate}

\author[1,2,3,$\ast$]{Michael P.H. Stumpf\ORCID{0000-0002-3577-1222}}

\address[1]{\orgdiv{School of BioSciences}, \orgname{University of Melbourne},
\orgaddress{\street{}, \postcode{3010}, \state{Victoria}, \country{Australia}}}
\address[2]{\orgdiv{School of Mathematics and Statistics},
\orgname{University of Melbourne},
\orgaddress{\postcode{3010}, \state{Victoria}, \country{Australia}}}
\address[3]{\orgname{Cell Bauhaus Pty Ltd},
\orgaddress{\state{Melbourne}, \country{Australia}}}

\corresp[$\ast$]{Corresponding author.
\href{mailto:mstumpf@theosysbio.org}{mstumpf@theosysbio.org}}

%\received{}{}{}
%\revised{}{}{}
%\accepted{}{}{}

% ---- Structured abstract ----------------------------------------------------
\abstract{
\textbf{Summary:}
Whether a dynamical model of a biological system can be understood as motion on
a potential landscape depends on its structural class. \texttt{FlowClass.jl} is a Julia package that
classifies a continuous-time system $\mathrm{d}\bm{x}/\mathrm{d}t =
\bm{F}(\bm{x})$ into a nested hierarchy of structural classes, namely gradient,
gradient-like, Morse-Smale, structurally stable, and general. It applies a
suite of tests, including Jacobian symmetry analysis, curl
quantification, fixed-point detection and stability classification, periodic-orbit
detection by Poincar\'e section with Floquet analysis, and stable and unstable
manifold computation, and returns the class together with the underlying
quantitative measures. The package is aimed at developmental and systems
biology, where it tests if a Waddington landscape
interpretation is justified for a given gene regulatory model and quantify the
non-gradient (curl) flux that distinguishes forward differentiation from
reprogramming.\\
\textbf{Availability and implementation:}
\texttt{FlowClass.jl} is open source under the MIT licence and is
available at \href{https://github.com/Theosysbio/FlowClass.jl}{https://github.com/Theosysbio/FlowClass.jl},
with an archived release on Zenodo. Documentation, worked examples, and a test
suite are included in the repository.\\
\textbf{Contact:} \href{mailto:mstumpf@theosysbio.org}{mstumpf@theosysbio.org}\\
%\textbf{Supplementary information:} Supplementary data are available at
%\textit{Bioinformatics} online.
}

\keywords{dynamical systems, Waddington landscape, cell fate, gradient systems,
Morse-Smale, Julia}

\maketitle

%=============================================================================
\section{Introduction}

Dynamical systems models of the form $\mathrm{d}\bm{x}/\mathrm{d}t =
\bm{F}(\bm{x})$ are ubiquitous across biology, physics, and engineering
\citep{Strogatz2015}. Dynamical systems that can be described in terms of a potential landscape are of particular interest.  
In developmental biology, following Waddington, they have become a popular and convenient descriptions of  Waddington's epigenetic landscape, which remains a dominant metaphor for cell differentiation
\citep{Waddington1957, Huang2007, Moris2016}. Much of the intuition built on
that metaphor implicitly assumes gradient dynamics, yet most gene regulatory
networks are not gradient systems: non-reciprocal regulatory interactions
generate curl, or flux, components that no potential can represent
\citep{Brackston2018, Wang2015}.

The distinction is not merely formal. For a gradient system the potential fully
determines the dynamics, and the minimum action paths for a transition and its
reverse coincide. Even for moderate curl fluxes, forward and reverse
paths separate, and differentiation and reprogramming follow
different routes through gene expression space \citep{Brackston2018,
Guillemin2021}. More generally, the structural class of a model  controls which
analytical tools can be applied: gradient and gradient-like systems admit global Lyapunov
functions and the machinery of catastrophe theory \citep{Rand2021}, whereas more
general systems do not \citep{Smale1967, PalisDeMelo1982}.

There is a lack of general-purpose tools  for classifying dynamical systems by their structural properties.
Existing Julia packages address related  problems, including solving ODEs and SDEs
in \texttt{DifferentialEquations.jl} \citep{Rackauckas2017}, chaos and
attractor analysis in \texttt{DynamicalSystems.jl} \citep{Datseris2018}, and
global set-oriented analysis in \texttt{GAIO.jl} \citep{GAIO2025}. None of these 
determine the qualitative structural nature of a specific vector field. \texttt{FlowClass.jl} fills this gap and places a system within a nested hierarchy of structural classes and reports the evidence behind the assignment.

\section{Features and implementation}

\texttt{FlowClass.jl} is written in  Julia \citep{Bezanson:2017,
Roesch2023} and depends only on \texttt{ForwardDiff.jl} for automatic
differentiation, \texttt{NLsolve.jl} for nonlinear solving, and
\texttt{OrdinaryDiffEq.jl} for trajectory integration. A system is defined by
wrapping a vector field in a \texttt{DynamicalSystem} object, after which a
single call to \texttt{classify\_system} runs the full analysis and returns a
\texttt{ClassificationResult}. The results contain the assigned class, 
detected invariant seta, the quantitative measures described below, and a
confidence level. The classification hierarchy runs from the most restrictive
class to the most general as gradient $\subset$ gradient-like $\subset$
Morse-Smale $\subset$ structurally stable $\subset$ general, and the pipeline
moves through the tests in sequence.

\paragraph{Jacobians}
 Jacobians are evaluated by automatic
differentiation, which avoids finite-difference error. A gradient system
$\bm{F} = -\nabla V$ has a Jacobian equal to the negative Hessian of $V$ and so
is symmetric everywhere. The scale-independent relative symmetry error $\|(J -
J^{\top})/2\|_F / \|J\|_F$ quantifies the departure from gradient structure and
is sampled across  regions of interest.

The package computes the curl directly in two and three dimensions, and uses the Frobenius norm of the antisymmetric part of
the Jacobian  in higher dimensions. The curl-to-gradient ratio measures the relative strength of rotational versus potential-driven
dynamics. This is key diagnostic for whether a landscape description is
appropriate.


 Fixed points are located by multi-start solving from Latin hypercube starting points within user-specified bounds, then
aggregated by clustering. Each point is classified from its Jacobian  as a stable or unstable node, saddle, focus, centre, or
non-hyperbolic point. A non-hyperbolic fixed point excludes the system from the
Morse-Smale class.

 Limit cycles are detected from recurrence along
integrated trajectories and refined on a Poincar\'e section, and their stability
is determined from the Floquet multipliers of the monodromy matrix. Gradient and
gradient-like systems cannot support periodic orbits, so their presence puts a
system at least in the Morse-Smale class.

For saddle points the package integrates the
stable and unstable manifolds from initial conditions perturbed along eigenvector
directions, detects homoclinic connections, and tests numerically whether the
manifolds intersect transversally, a requirement for Morse-Smale structure.

These analyses are combined by \texttt{classify\_system}, with lighter-weight
\texttt{quick\_classify} and \texttt{get\_system\_class} functions providing more limited but  rapid
screening. Queries such as \texttt{is\_gradient},
\texttt{allows\_periodic\_orbits}, and \texttt{has\_landscape\_representation}
translate the result into the landscape interpretation, and
\texttt{print\_classification} produces a formatted report.

\section{Example: stem cell differentiation}

To illustrate the package on a biological model, \texttt{FlowClass.jl} ships with
the embryonic stem cell differentiation model of \citet{Brackston2018}, itself
built on \citet{Chickarmane2012}, which describes the pluripotency factors Nanog,
Oct4-Sox2, and Fgf4  and the differentiation marker Gata6. The four-dimensional system has multiple stable fixed
points corresponding to pluripotent and differentiated cell fates, separated by a
saddle that represents the transition state;  the locations of these states
shift as  external signals change.

Running \texttt{classify\_system} on this model recovers the bistable structure
and  returns a non-negligible curl-to-gradient ratio. The practical consequence is that the landscape, while a useful
organizing picture, does not by itself fix the dynamics, and the paths taken
during differentiation and during reprogramming are distinct
\citep{Brackston2018, Guillemin2021, Vittadello2025}. Being able to quantify
this flux directly from the model, rather than assuming it, is what makes
landscape-based inference from single-cell data well posed \citep{Liu2024}.

\begin{figure}[t]
\centering
\includegraphics[width=\linewidth]{JOSS-Fig1.pdf}
\caption{Waddington's epigenetic landscape with potential wells corresponding to
cell fates. On the left the flow is purely gradient and is fully determined by
the landscape. On the right the flow carries a curl component, so the dynamics
are no longer fixed by the landscape and forward and reverse transition paths
differ. \texttt{FlowClass.jl} distinguishes these regimes automatically and
quantifies the curl contribution.}
\label{fig:landscape}
\end{figure}

\section{Discussion}

The qualitative statements we can make about a biological model depend
critically on its structural class. \texttt{FlowClass.jl} makes that class an
explicit, testable property rather than an unstated assumption, and reports the
quantitative measures, Jacobian symmetry, curl ratio, fixed-point types, and
orbit stability, on which the assignment rests. Where existing tools ask how a
system behaves, \texttt{FlowClass.jl} asks what kind of system it is, and the two
are complementary: a structural classification is a natural first step before
bifurcation continuation, landscape reconstruction, or action-path calculation.
For developmental biology in particular, the package offers a direct check on
when the Waddington metaphor is sound and a measurement of the non-equilibrium
flux that it omits.

\section*{Acknowledgements}

The author thanks members of the Theoretical Systems Biology group for testing
and feedback.

\section*{Author contributions}

M.P.H.S. conceived the work, implemented the software, and wrote the manuscript.

\section*{Conflict of interest}

The author is a director, shareholder and chief scientific officer  of Cell Bauhaus Pty Ltd. No other conflicts are
declared.

\section*{Funding}

This work was supported by the Australian Research Council through an ARC
Laureate Fellowship (FL220100005).

\section*{Data availability}

No new data were generated in support of this work. The software, documentation,
worked examples, and test suite are available at
\url{https://github.com/Theosysbio/FlowClass.jl} and in an archived Zenodo
release (10.5281/zenodo.20807340).

\bibliographystyle{abbrvnat}
\bibliography{paper}

\end{document}
